%% chapters/21-spectral-clustering.tex
\chapter{Spectral Clustering: Three Perspectives}
\label{ch:specclust}

\whereWeAre{%
The destination. \emph{Spectral clustering} is what happens when you take the spiked-covariance picture seriously: the signal eigenvectors of a similarity matrix, graph Laplacian, or kernel reveal cluster structure in the data. We follow von Luxburg and present the algorithm three different ways, so that if one explanation doesn't click, the other two will.%
}

\PLACEHOLDER{New chapter. Primary source: von Luxburg's \emph{A Tutorial on Spectral Clustering} (the canonical pedagogical reference).}

\section{The algorithm in one paragraph}
\PLACEHOLDER{Build similarity graph; compute graph Laplacian; take its smallest $k$ eigenvectors; cluster the $n$ embedded points with $k$-means.}

\section{Perspective 1: graph cut}
\PLACEHOLDER{Normalized cut interpretation; spectral relaxation of an NP-hard problem.}

\section{Perspective 2: random walk}
\PLACEHOLDER{Random walk on the similarity graph; the Laplacian eigenvectors are stationary modes.}

\section{Perspective 3: perturbation theory}
\PLACEHOLDER{The Laplacian of a block-structured graph has a specific spectrum; perturbations don't break the picture too much. This is where the BBP intuition feeds in directly.}

\section{Why this is the destination of the entire book}
\PLACEHOLDER{The full chain: Gaussian foundations $\to$ random matrices $\to$ spiked covariance $\to$ BBP detectability $\to$ clusters revealed by top eigenvectors. Spectral clustering is exactly the algorithmic payoff.}

\section{Simulation: spectral clustering on a synthetic dataset}
\PLACEHOLDER{Colab box: two-blob dataset, build kernel, compute Laplacian eigenvectors, $k$-means in spectral coordinates.}

\section{Brief warm-ups}
\PLACEHOLDER{3--5 short questions.}

\chapterRecap{%
\begin{itemize}
  \item \PLACEHOLDER{Spectral clustering = top/bottom Laplacian eigenvectors + $k$-means.}
  \item \PLACEHOLDER{Three independent reasons it works.}
  \item \PLACEHOLDER{Random matrix theory is the explanation for \emph{when} it works.}
\end{itemize}%
}

\section*{Exercises}
\PLACEHOLDER{8--15 longer exercises.}
